\begin{statement}[C.6.4; Proved: the maximum as a height inequality]
For $t\in K\setminus\{0,1\}$,
\begin{equation}\label{eq:ch-bound}
 \widehat h_O(P_t)\geq
 \frac1{[K:\Q]}\sum_{v\nmid\infty}n_v\lambda_v(P_t)
                    -\frac12\log2.
\end{equation}
The inequality is strict on this smooth base. Over $\Q$ it
is $\widehat h_O(P_t)\geq\log A_t-\frac12\log2$.
The generic height is $\widehat h_O(P)=1/12$, and the
variation divisor is $D_P=(\infty)/12$. Consequently
\[
 \widehat h_O(P_t)=\frac1{12}h(t)+O(1)
 \qquad(t\in\overline\Q\setminus\{0,1\}),
\]
where the constant is uniform over algebraic $t$ and is for
the fixed family and section.
\end{statement}
\begin{proof}
Apply the analytic maximum theorem C.8.5 at every embedding
and substitute into \eqref{eq:ch-global}. Equality would
require $j(\sigma(t))=1728$ for every embedding, hence $t=1$,
which has been excluded.

For the generic local heights at $0$ and $1$, the minimal
discriminant orders are $8$ and $3$, and the section meets
the nonidentity component of types $IV^*$ and $III$.
Their local corrections are $2/3$ and $1/4$, respectively,
so their normalized local heights are zero. At infinity,
put $s=1/t$, $X=s^{-2}x$, $Y=s^{-3}y$. The new equation is
\[
 y^2=x^3-\tfrac34(1-s)x+\tfrac14(1-s)^2,
 \qquad P=(0,(1-s)/2),\qquad\Delta=27s(1-s)^3.
\]
At $s=0$ the node is $(1/2,0)$, whereas $P=(0,1/2)$ is
nonsingular. Its height contribution is $1/12$.
At all other function-field places the point is integral
and the discriminant is a unit. This proves
$D_P=(\infty)/12$ and the generic height.

Theorem 1.1 of \citet{demarco2017} applies with $K=\Q$,
$B=\mathbb P^1$, the displayed elliptic surface and rational
section, since the generic height is positive. It equips
this divisor with continuous adelic semipositive metrics
whose associated height is exactly $\widehat h_O(P_t)$.
Comparing with the usual height for $(\infty)/12$ gives the
bounded difference. This is the specialization here of the
variation theorem of Tate, with the local refinements of
Silverman reproduced as Theorems 2.2 and 2.3 in that paper:
outside finitely many places the local correction is zero,
and at each remaining place it extends continuously after
subtracting its prescribed logarithmic singularity.
Our inequality is an explicit archimedean estimate with
constant $\frac12\log2$; the general bounded-variation
statement alone does not supply that constant.

We spell out the comparison with the opened primary author sources.
Theorem 4.1 of \citet[pp.~183--184]{callSilverman1993} concerns a
smooth projective base curve $T$, a degree-one base height $h_T$,
a section, and a fibrewise morphism whose pullback multiplies the
chosen divisor class by a real number $\alpha>1$. It gives
\[
 \lim_{h_T(t)\to\infty}
 \frac{\widehat h_t(P_t)}{h_T(t)}
       =\widehat h_{\Q(T)}(P).
\]
Here take $T=\mathbb P^1$, the class of $(O)$, and the generic
morphism $[2]$, with $[2]^*(O)\sim4(O)$ and $\alpha=4$.
The smooth fibres are precisely the complement of $0,1,\infty$;
the section extends on the regular projective elliptic surface.
All the hypotheses hold, and the limit is $1/12$.
The unnumbered remark immediately following that theorem states
Tate's stronger elliptic-family formula
$\widehat h_t(P_t)=h_{T,D_P}(t)+O_P(1)$, with
$\deg D_P=\widehat h_{\Q(T)}(P)$. This is an author-verified
reproduction of Tate's result, not a claim to have read the
inaccessible original scan \citep{tate1983variation}.
The same distinction applies to \citet{silverman1983specialization}:
the limit theorem is reproduced by the author on slide~22 of
\citet{silverman2006heights}, while slide~21 states the uniform
canonical-versus-Weil-height bound in terms of the base height.
For our explicit bounded remainder and local normalization, the
fully opened Theorems~1.1, 2.2 and 2.3 of
\citet{demarco2017}, applied above, remain the cited proof.
\end{proof}
